%% =====================================================================
%%  T-04-04  The Self-Audit
%%  -------------------------------------------------------------------
%%  One idea per file. Core is mandatory. Slots in canonical order.
%%  Max 700 words. No layout commands. No filler. New source material
%%  for this entry goes in sources/<BOOK>/T-04-04.tex -- never by
%%  rewriting this file (docs/RULES.md R0.1).
%% =====================================================================
%% @kind: topic
%% @id: T-04-04
%% @title: The Self-Audit
%% @subtitle: Find your own load-bearing points before someone else does
%% @chapter: c04
%% @order: 40
%% @status: stable
%% @sources: B1, B3
%% @updated: 2026-09-19
%% @allow-rewrite: slot migration to the seven-question format (docs/RULES.md section 2)

\topic{T-04-04}{The Self-Audit}{Find your own load-bearing points before someone else does}


\begin{core}
The defensive half of this chapter is a short list of your own weak points, written down while you are calm. It is the highest-yield item in Part I, because every technique aimed at you must pass through one of them.
\end{core}

\begin{mechanism}
Techniques do not work on people generically; they work on dispositions specifically. A person who knows that they cannot bear to be thought unfair will notice immediately when a counterpart is building a case around fairness, and that notice converts an automatic reaction into a choice. The audit does not remove the weakness --- it cannot --- it moves it from the automatic system to the deliberative one, which is where the counters in Part VI operate.
\end{mechanism}

\begin{application}
  \item List the subjects that change your register. Use the last three months of real conversations, not a theory about yourself.
  \item For each, write the request that would be hardest to refuse and the accusation that would be hardest to answer.
  \item Write the pre-commitment: what you will do the next time one of these is loaded. Usually it is delay.
  \item Give the two or three most serious items air with one person who has no use for them.
  \item Review the list when circumstances change --- a new job, a new relationship, a loss --- because the points move.
\end{application}

\begin{feedback}
  \item You can name the accusation that would most upset you.
  \item You know which request you find hardest to refuse.
  \item You can identify the two or three people whose opinion of you actually changes your behaviour.
  \item You know what you are unwilling to have examined in public.
\end{feedback}

\begin{failure}
Self-knowledge is unreliable in exactly the areas the audit targets, which is why the list tends to be incomplete and why the third-party check above is not optional. There is also a failure mode in over-auditing: a person who monitors every reaction becomes unable to act, and hesitation is more exploitable than any single weakness. The list should be short --- five items, not fifty --- and read before difficult exchanges rather than during them.
\end{failure}

\begin{countermeasures}
  \item If you are the target, the audit is the counter. Most of Part VI assumes you have done it.
  \item Keep a standing delay rule for any request that touches a listed point. The rule does the refusing so you do not have to.
  \item Ask a person who has known you a long time and gains nothing to name your most obvious one. External observers see these faster than their owners do.
  \item Treat sudden intense interest in one of your listed points as a boundary condition: slow down, get it in writing, involve a third party.
\end{countermeasures}

\sources{B1, B3}
\seesources
\seealso{T-04-01, T-23-01, T-24-01}
